\scanpage{191}
\mspage{140}
Regardant la restriction des deux morphismes de diagrammes $(\ast\ast\ast)$ aux parties des \inserted{trois} diagrammes de type $\Gamma^0$ au dessus de la queue $\{3,4\}$ de $\Gamma$, on trouve le diagramme à traits pleins que voici :
\diagram{\begin{tikzcd}[ampersand replacement=\&,row sep=small,column sep=small]
\&\int\Sigma X_\ast\arrow[dl,"W"']\arrow[dr,"W"]\arrow[dd,"{\varphi_{\Sigma Q}}"']\&\&\Sigma\widetilde{\mathscr X}\arrow[dl,"W"',postaction={decorate,decoration={snake,amplitude=.5pt,segment length=3pt}}]\arrow[dr,dashed,"W"]\arrow[dd,dashed,"{\Sigma\widetilde\varphi_Q}"]\&\\
\mathscr X\arrow[dd,"{\unsure{\varphi|\rho}_Q}"']\&\&\widetilde{\mathscr X}^0\arrow[dd,"{\varphi_{Q^0}=(\widetilde\varphi_Q)^0}"]\&\&\widetilde{\mathscr X}\arrow[dd,densely dotted,preaction={draw,solid,thin},"{\widetilde\varphi_Q}"]\\
\&\Sigma Y\arrow[dl,"W"']\arrow[dr,"W"]\&\&\Sigma Y\arrow[dl,dashed,"W"']\arrow[dr,dashed,"W"]\&\\
Y\&\&Y^0\&\&Y\rlap{,}
\end{tikzcd}}
que j'ai complété par un deuxième volet en pointillés. Dans ce diagramme de quatre carrés accolés, les côtés horizontaux des carrés \deleted{\unsure{représentent}} des flèches de $\mathit{Cat}$ qui sont \unsure{des|dans} $W$. On trouve ainsi :
\[
\varphi_Q\in W\Longleftrightarrow\varphi_{\Sigma Q}\in W\Longleftrightarrow
(\varphi_{Q^0}=\widetilde\varphi_Q^0)\in W\Longleftrightarrow
\Sigma\widetilde\varphi_Q\in W\Longleftrightarrow\widetilde\varphi_Q\in W .
\]
On trouve aussi que les cinq flèches verticales sont isomorphes, en tant que flèches de $\mathrm{Hot}_W$.

Pour résumer :

\uline{Proposition 14} Considérons un \inserted{diagramme} \unsure{(carré commutatif)} 
\diagram{\begin{tikzcd}[ampersand replacement=\&,column sep=small,row sep=tiny]
\&\&X_0\arrow[dl]\arrow[dr]\&\\
Q\deleted{\unsure{Q\to}}={}\&X_1\arrow[dr]\&\&X_2\arrow[dl]\\
\&\&Y\&
\end{tikzcd}}
dans $\mathit{Cat}$, et le diagramme \inserted{\textquotedblleft opposé\textquotedblright} déduit par $X\mapsto X^0$
\diagram{\begin{tikzcd}[ampersand replacement=\&,column sep=small,row sep=tiny]
\&\&X_0^0\arrow[dl]\arrow[dr]\&\\
Q^0={}\&X_1^0\arrow[dr]\&\&X_2^0\arrow[dl]\\
\&\&Y^0\&{,}
\end{tikzcd}}
% archive sheet number: 191
